% New post-checkpoint-13 scalar-cover extension. Independent audit requested.
\begin{proposition}[Proved under stated hypotheses: finite matrix groups with bounded scalar kernels]
\label{D:bounded-scalar-projective-growth}
Let $Q$ and $p$ satisfy the two hypotheses of
Proposition~\ref{D:bounded-centre-projective-growth}: constant
prime-to-$p$ second cohomology is finite, and every open normal
subgroup has a quadratic finite-field representation bound.
Fix $B\ge1$. For $q=p^f$, let
\[
 \widetilde S\subset\operatorname{GL}_r(\mathbf F_q),\qquad
 S=\operatorname{im}\bigl(\widetilde S\longrightarrow
                         \operatorname{PGL}_r(\mathbf F_q)\bigr).
\]
Suppose that the scalar kernel of $\widetilde S\to S$ has order at
most $B$. Then, with a constant $C_B$ independent of $r,f$ and
$\widetilde S$,
\[
 \#\bigl\{\operatorname{Hom}_{\rm cont}(Q,S)
             \text{ up to conjugacy in }
                    \operatorname{PGL}_r(\mathbf F_q)\bigr\}
       \le q^{C_Br^2}.                                  \tag{D.SC.1}
\]
We mean the restriction of the ambient conjugacy equivalence
relation to these homomorphisms. We do not require $S$ to be normal
in $\operatorname{PGL}_r(\mathbf F_q)$.

Consequently, fix also $\epsilon>0$ and consider any family of
nonabelian finite simple groups $S$ with such realizations satisfying
$|S|\ge q^{\epsilon r^2}$. The quotient of $Q$ given by all distinct
simple quotient kernels of these types is topologically finitely
generated.
\end{proposition}
\begin{proof}
The scalar kernel is a subgroup of the cyclic group $\mathbf F_q^\times$.
It is therefore central and cyclic of prime-to-$p$ order $d\le B$.
The proof of Proposition~\ref{D:bounded-centre-projective-growth}
constructs one open normal subgroup $H_B\subset Q$ on which every
constant degree-two class with cyclic coefficients of such an order
vanishes. The pullback of
\[
 1\longrightarrow\mathbf Z/d\mathbf Z\longrightarrow
        \widetilde S\longrightarrow S\longrightarrow1
\]
by a homomorphism of $Q$ thus splits continuously over $H_B$.
Every restricted projective representation lifts to a linear
representation of $H_B$ over the original field $\mathbf F_q$.
Its ambient projective conjugacy class is consequently among at
most $q^{c_{H_B}r^2}$ possibilities.

Choose a representative of each such ambient restricted class.
Conjugating a homomorphism of $Q$ so that its restriction equals
this representative may move its full image out of the originally
specified subgroup $S$. This causes no loss in an upper bound:
we count all its possible extensions as homomorphisms into the
ambient projective linear group. If $t=[Q:H_B]$, their values on
$t-1$ fixed nonidentity coset representatives give at most
$|\operatorname{PGL}_r(\mathbf F_q)|^{t-1}$ choices. Thus (D.SC.1)
holds with $C_B=c_{H_B}+t-1$.

Ambient-conjugate epimorphisms have the same kernel, so the number
of distinct kernels for a simple group in the family is at most
$|S|^{C_B/\epsilon}$. The generating-tuple estimate used in
Proposition~\ref{D:bounded-matrix-simple-product} gives a uniform
number of generators for all finite products with these
multiplicities; its precise primary locator is the proof of
Proposition~2.6(ii), pp.~12--13, of
\citet{LucchiniMarionTracey2020D}. Goursat's lemma identifies the
image associated with distinct simple quotient kernels with the
full product on every finite subcollection. Compactness then gives
a generating tuple for the profinite product.
\end{proof}

\begin{corollary}[Proved: unitary quotients with bounded natural centre]
\label{D:bounded-centre-unitary-simple-product}
Let $Q$ satisfy the hypotheses of
Proposition~\ref{D:bounded-scalar-projective-growth}, and fix $B$.
For all $r\ge2$, $q=p^f$, and $\gcd(r,q+1)\le B$, the number of
homomorphisms
\[
           Q\longrightarrow\operatorname{PSU}_r(q)
\]
up to conjugacy in $\operatorname{PGL}_r(\mathbf F_{q^2})$ is
at most $q^{C'_Br^2}$. The quotient obtained by taking all distinct
nonabelian simple quotients of these unitary types is
topologically finitely generated, with no bound on $r$.

In particular these assertions hold for the original-boundary tame
group of a normal proper surface in characteristic $p$, and for
the fundamental group of a proper scheme of dimension at most two
after removing a finite set of codimension at least two in every
irreducible component it meets.
\end{corollary}
\begin{proof}
The natural special unitary group is a subgroup of
$\operatorname{GL}_r(\mathbf F_{q^2})$. Its scalar kernel in the
projective linear group has order $d=\gcd(r,q+1)$: a scalar
$\lambda$ preserves the hermitian form precisely when
$\lambda^{q+1}=1$, and has determinant one precisely when
$\lambda^r=1$. Its projective image is $\operatorname{PSU}_r(q)$.
These statements, and the order formula
\[
 |\operatorname{PSU}_r(q)|
   =\frac{q^{r^2-1}}{d}
          \prod_{i=2}^r\bigl(1-(-1)^i q^{-i}\bigr),      \tag{D.SC.2}
\]
are proved in \citet[Chapter~10, p.~118]{Taylor1992ClassicalB}.
The same source, Theorem~10.20, pp.~127--128, proves simplicity
except for $(r,q)=(2,2),(2,3),(3,2)$; we include only the
nonabelian simple members in the product assertion.

Applying the proposition over the field of size $q^2$ gives
the representation bound $q^{2C_Br^2}$. For completeness its
order comparison is uniform in rank. Odd-index factors in
(D.SC.2) are at least one, and
\[
 \prod_{\substack{2\le i\le r\\i\text{ even}}}(1-q^{-i})
       \ge 1-\sum_{j\ge1}2^{-2j}=\frac23.
\]
Since $d\le r$, $2r\le2^r\le q^r$, and
$r^2-r-1\ge r^2/4$ for $r\ge2$, we obtain
\[
           |\operatorname{PSU}_r(q)|\ge q^{r^2/4}
                         =(q^2)^{r^2/8}.
\]
The simple-product assertion follows with $\epsilon=1/8$.

For a normal proper surface with an original tame boundary,
Theorem~\ref{D:surface-tame-cross-characteristic-H2} supplies
constant second-cohomology finiteness. Each open normal subgroup
is a quotient of the original-boundary tame group of its finite
normal proper covering, exactly as in
Corollary~\ref{D:surface-bounded-centre-simple-product}.
Theorem~\ref{D:tame-global-quadratic-representation-bound}
supplies the required hereditary quadratic counts.

For a possibly nonnormal proper scheme of dimension at most two,
we repeat the finite proper normalization descent of
Corollary~\ref{D:nonnormal-bounded-centre-projective-growth},
now for torsors under $\operatorname{PGL}_r(\mathbf F_{q^2})$.
On normal two-dimensional components the just-proved bound
applies. Lower-dimensional components are proper, so their fixed
finite generator bounds give $q^{O(r^2)}$ classes. For any fixed
tuple of restricted torsors, descent isomorphisms on the fixed
double product have at most
$|\operatorname{PGL}_r(\mathbf F_{q^2})|^c\le q^{2cr^2}$
possibilities, where $c$ is its number of connected components.
Triple-overlap relations only reduce this number. Effective
finite proper descent \citep[IX, Theorem~4.7]{C_SGA1} and the
same order comparison finish the proof. No second-cohomology
assertion for the nonnormal total scheme is used.
\end{proof}

\begin{remark}[Exact central scope]
The constants may depend on the fixed bound $B$. The unitary
statement does not cover an unbounded sequence of values
$\gcd(r,p^f+1)$. The scalar-cover proposition also does not count
arbitrary monolithic nonabelian crowns; its product conclusion
concerns simple quotients only.
\end{remark}
